\chapter{Threat model and security requirements}
\label{chap:threat-model}

\section{System model}

The operator supplies a trusted task goal \(G\). At agent step \(t\), the
browser produces an observation \(O_t\), the model and scaffold produce a
proposed action \(a_t\), and the privileged executor may update environment
state \(S_t\):

\[
  O_t = \mathsf{observe}(S_t), \qquad
  a_t = \mathsf{agent}(G,O_{\leq t},A_{<t}), \qquad
  S_{t+1} = \mathsf{execute}(S_t,a_t).
\]

In an ungoverned loop, the last transition follows directly from the model's
proposal. \systemname\ inserts an observation transformation
\(\mathsf{guard}_{O}\) and an action decision
\(\mathsf{guard}_{A}\):

\[
  \widehat{O}_t = \mathsf{guard}_{O}(O_t,P,C_t), \qquad
  d_t = \mathsf{guard}_{A}(a_t,S_t,P,C_t),
\]

where \(P\) is live policy and \(C_t\) is run-local governance state, including
taint, approvals, budgets, and cancellation. The executor is called only when
\(d_t\) yields \emph{allow} or a still-valid approved scope. Block, timeout,
cancellation, unavailable required context, and unresolved authority do not
execute \(a_t\).

This model separates \emph{proposal} from \emph{authority}. The model can
propose any syntactically valid tool call, including one influenced by
malicious content. The proposal has no effect until a component outside the
model evaluates it. The separation is the primary architectural security
property; the quality of individual detector and risk rules is a secondary,
empirical property.

\section{Actors, assets, and trust boundaries}

\begin{table}[htbp]
\centering
\caption{Actors and trust assumptions.}
\label{tab:trust-boundaries}
\small
\begin{tabularx}{\textwidth}{@{}L{0.25\textwidth}L{0.18\textwidth}Y@{}}
\toprule
Component or actor & Trust & Rationale \\
\midrule
Operator request and live steering & Trusted authority &
Defines the task and may provide new instructions, but cannot override a
configured hard block merely by approving a pending action. \\

Local policy and \systemname\ middleware & Trusted computing base &
Defines the allowed authority and mediates wrapped model inputs and
model-proposed privileged tool actions at the governed boundary. It does not
mediate direct human page interaction or browser effects outside those
channels. \\

Browser-agent model and scaffold & Untrusted decision maker &
May hallucinate, follow injected content, retry denials, or produce malformed
actions. \\

DOM, accessibility text, title, URL-carried content, extracted text, downloads,
and tool results & Untrusted data &
May contain first-party, third-party, user-generated, or attacker-controlled
content, including stored attacks. \\

Screenshot pixels & Excluded from ordinary model and operator-telemetry paths &
DOM sanitisation cannot clean text rendered into pixels. Optional local trace
retention does not make a frame trusted. \\

Browser and tool implementation & Privileged executor &
Performs effects after mediation and therefore remains inside the trusted
computing base. \\

Extension side panel & Operator interface, not an authority source &
Displays redacted decisions and submits responses that the server validates;
it cannot mint broader consent locally. \\

Trace consumers & Honest but curious &
May inspect records but should not receive raw secrets, injected payloads,
screenshots, or model rationale by default. \\

Other same-user processes & Outside the accepted boundary &
Loopback binding is accepted for this local research prototype; it does not
isolate the service from a compromised process with the same operating-system
account. \\
\bottomrule
\end{tabularx}
\end{table}

Protected assets include the operator's task authority, authenticated browser
session, private or security-sensitive data, local files available for upload,
integrity of remote state, consent decisions, and the confidentiality of audit
records. Availability is a secondary asset: a defence that blocks every task
is not useful, but fail-closed denial is preferred to an unauthorised side
effect when required security context is missing.

\section{Adversary}

The primary adversary controls content that the agent may encounter during a
legitimate task. This includes a page or part of a page, a user-generated
record, a link label, a tool result, or content later returned by extraction.
The adversary knows the general agent design and may write natural-language
instructions intended to:

\begin{itemize}
  \item replace or extend the operator's goal;
  \item impersonate system, developer, tool, or user authority;
  \item demand secrecy or suppression of warnings;
  \item request navigation, disclosure, code execution, file access, or a
  consequential same-origin action;
  \item persist influence through history or a second model call; or
  \item exploit ambiguities in page structure and target semantics.
\end{itemize}

The attacker need not compromise the model provider, \systemname\ process,
browser binary, operating system, or local policy file. Those stronger
attackers are out of scope. The benchmark adversary is also bounded by each
suite's instantiated attack goal and environment; the evaluation does not
claim coverage of every adaptive or obfuscated payload.

\section{Security objectives}
\label{sec:security-objectives}

The design is governed by eight objectives.

\begin{enumerate}[label=\textbf{O\arabic*:},leftmargin=16mm]
  \item \textbf{Complete action mediation.} Every model-proposed privileged
  tool action entering \texttt{GovernedTools} is evaluated immediately before
  execution. There is no model-accessible parallel path to the wrapped tool
  through the implemented agent scaffold.

  \item \textbf{Complete terminal audit.} Every such governed tool attempt has
  exactly one privacy-preserving attempted event and one terminal event,
  including blocks, denials, timeouts, exceptions, and cancellation.

  \item \textbf{Source does not imply authority.} Browser, file, extraction,
  and tool-result content remains untrusted even when it is necessary for the
  task or originates on the current first-party domain.

  \item \textbf{Controlled sensitive-data flow.} Recognised sensitive values
  cannot enter an untrusted sink solely because a page or the model requested
  it. In enforce mode, credentials, one-time codes, access tokens, payment
  values, private keys, and protected placeholders are completed by the human
  directly, not approved for agent entry.

  \item \textbf{Fail-safe authority.} Consent is explicit, narrow, expiring,
  run-local, and non-transferable. Silence, timeout, malformed responses,
  shutdown, and a stop race do not execute the pending action.

  \item \textbf{Observation post-condition.} Sanitisation counts as successful
  only if at least one attributable untrusted region changed and re-detection
  over the resulting model input is clear. A known residual or unattributable
  current browser observation is withheld before the model call.

  \item \textbf{Monitoring confidentiality.} Dedicated enforcement fields
  retain redacted summaries, finding metadata, and run-keyed digests, not raw
  detector payloads, recognised secrets, or unredacted parameters. Model
  metadata omits HTTP(S) queries and fragments; telemetry may retain benign ones
  after sensitive-value masking. Both omit non-web suffixes. Recursive
  redaction covers general events and history. Arbitrary model echoes and the
  local process log are outside this invariant.

  \item \textbf{Utility-aware enforcement.} Read-only inspection remains
  available where policy permits it, and ambiguous risk is not converted into
  a blanket block. The evaluation must report legitimate utility beside
  security outcomes.
\end{enumerate}

Objectives O1--O7 are implementation invariants within the channel scopes
defined above. O8 is a design and evaluation criterion rather than a universal
guarantee: no runtime can ensure task completion for an arbitrary model,
website, and policy.

\section{Taint and capability model}

\systemname\ assigns labels to sources and action parameters. A web observation
introduces \texttt{untrusted\_web}. A positive detector finding adds
\texttt{prompt\_injection\_suspected} and is retained per origin so that a
later clean viewport does not erase possible influence on model history.
Recognised values add categories such as \texttt{sensitive:credential},
\texttt{sensitive:one\_time\_code}, \texttt{sensitive:access\_token},
\texttt{sensitive:payment\_data}, or \texttt{sensitive:private\_key}. Local
uploads add \texttt{local\_file}.

Actions expose capabilities rather than being judged only by tool name.
Relevant capabilities include read-only inspection, navigation, page write,
data disclosure, local-file read or write, code execution, and task
completion. A navigation action also exposes a destination origin; an indexed
page action exposes bounded target semantics such as element role, label,
input type, and consequential button language.

This is dynamic taint tracking in an engineering sense, not a proof of
language-level non-interference. Model memory can transform data, a detector
can miss a source, and target semantics can be ambiguous. The labels support
explicit rules such as ``injection-tainted cross-origin disclosure requires
intervention'' and ``recognised credentials require takeover'' without
asserting that all semantic flows are captured.

\section{Threat--control mapping}

\begin{table}[htbp]
\centering
\caption{Principal threats and intended controls.}
\label{tab:threat-controls}
\small
\begin{tabularx}{\textwidth}{@{}L{0.31\textwidth}Y@{}}
\toprule
Threat & Control strategy \\
\midrule
Indirect injection in a current DOM/read observation &
Untrusted-region extraction, deterministic detection, attributable
sanitisation, clean re-detection, and origin-scoped taint. \\

Partial or unsafe sanitisation &
Withhold the current observation, open a persistent security incident, and
obtain a fresh observation after operator control; remove an unresolved
bounded extraction/history value in full. \\

Stored or second-order tool output &
Wrap all tool-result text in untrusted protocol regions and apply the same
boundary on later model calls. \\

Image or canvas instruction &
Exclude screenshots and image-bearing tool results from the ordinary model
path; declare visual attacks outside evaluated coverage. \\

Secret or token disclosure &
Recursive value/key classification, DOM field semantics, guarded free-text
keystrokes and page-code semantics, bounded local-text upload inspection,
direct human takeover, and recursive redaction. \\

Injected cross-origin or privileged action &
Sticky taint, capability classification, live target/origin analysis, dynamic
block or approval, and static domain rules. \\

Redirect, popup, or link-based domain bypass &
Mirror domain policy into the browser navigation watchdog in addition to
pre-tool evaluation. \\

Approval replay or page change &
Bind grants to run, action, origin/path, target fingerprint, capabilities,
taints, categories, and selected parameter digests; re-evaluate immediately
before execution. \\

Stop, timeout, or preflight exception &
Emit attempted audit before fallible work; fail closed and emit one terminal
outcome without invoking the tool. \\

Leak through observability &
Redact recursively, suppress raw page titles and non-web URL suffixes, use
run-keyed digests, and omit screenshots and free-form model rationale by
default. \\
\bottomrule
\end{tabularx}
\end{table}

\section{Assumptions and explicit non-claims}
\label{sec:nonclaims}

The following limits are part of the security statement rather than defects
hidden from it.

\begin{itemize}
  \item \textbf{Detection is incomplete.} Signatures cannot prove an
  observation benign. Novel, obfuscated, task-relevant, or adaptive attacks
  can be missed.

  \item \textbf{Visual input is not evaluated.} Screenshot and canvas attacks
  are outside the ordinary model path and the reported campaigns use
  DOM/accessibility observations. Disabling vision can also remove legitimate
  information.

  \item \textbf{Data-loss prevention is bounded.} The implementation recognises
  specified values, keys, field semantics, action types, filenames, and at most
  256 KiB of authorised local UTF-8 upload content. A semantically private
  value with no recognised marker may be missed.

  \item \textbf{Takeover is not post-handoff non-interference.} It prevents the
  intercepted model-proposed action from entering a protected value. An
  operator note is ordinary model context, and a non-password value entered
  directly into the page can reappear in a later DOM or accessibility
  observation. The system therefore makes no general claim that human-entered
  data remains outside all subsequent model context.

  \item \textbf{The executor is trusted.} Once an action is allowed, a remote
  site or browser handler may behave differently from its visible label.
  \systemname\ does not provide transactional isolation or verify server-side
  implementation.

  \item \textbf{Localhost is an accepted prototype boundary.} Origin checks and
  input validation protect the intended extension/API interaction, but a
  compromised same-user process can call the loopback service or modify local
  files. This is not an operating-system sandbox.

  \item \textbf{Policy does not fully infer intent.} Generic browser controls
  have imperfect business semantics. Operators or deployments must configure
  sensitive domains, trusted destinations, and prohibited capabilities
  appropriate to their tasks.

  \item \textbf{No universal or human-factors claim follows.} The experiments
  cover specific models, seeds, suites, policies, and scaffolds. No participant
  study measures trust, workload, comprehension, fatigue, usability, or
  rubber-stamping.
\end{itemize}

These boundaries govern \cref{chap:methodology}: AgentDojo tests controlled
tool mechanisms and WASP tests browser-native transfer; neither validates the
excluded visual or same-user compromise scenarios.
